\chapter{Moduli diagnostici e superficie minima}
\label{cap:16-moduli-diagnostici}

Questo capitolo chiude la Parte III con due argomenti che il resto non ha
collocato: il modulo pi\`u piccolo del progetto, e la configurazione che non ha
moduli affatto.

\section{\code{build-metadata}}

Quarantatr\'e righe in tre classi. Espone un endpoint,
\code{/modules/build-metadata}, che riporta l'identit\`a della build in
esecuzione.

La sua utilit\`a \`e proporzionale alla combinatoria che il sistema modulare
introduce. Un binario di \nf{} \`e definito da: la versione del codice, la
selezione dei moduli, il tipo di build (JVM o immagine nativa), il livello di
ottimizzazione, il garbage collector. Il \cref{cap:24-footprint} confronta
configurazioni che differiscono \emph{solo} per l'ultimo di questi assi. In un
tale contesto, l'assenza di un modo diretto per interrogare un'istanza in
esecuzione sulla propria identit\`a \`e una fonte reale di errore sperimentale:
il rischio non \`e teorico, ed \`e quello di attribuire a una politica una
differenza dovuta al fatto che una delle due istanze non era quella che si
credeva.

Che sia un modulo, e non un endpoint del nucleo, \`e coerente: la
diagnostica sulla build serve a chi conduce esperimenti, non alla piattaforma.
Un deployment di produzione pu\`o non compilarla.

\section{La configurazione senza moduli}

\subsection{Che cosa resta}

Compilando con \code{-PcontrolPlaneModules=none} si ottiene una piattaforma con
le seguenti propriet\`a:

\begin{center}
\small
\begin{adjustbox}{max width=\textwidth}
\begin{tabular}{@{}lp{8.4cm}@{}}
\toprule
\textbf{Capacit\`a} & \textbf{Stato} \\
\midrule
Registrazione e ciclo di vita & Completa. \\
Invocazione sincrona & Completa; dispatch diretto, senza coda n\'e ammissione. \\
Invocazione asincrona & \code{501 Not Implemented}. \\
Idempotenza, retry, timeout & Complete. \\
Modalit\`a \code{LOCAL} ed \code{EXTERNAL} & Complete. \\
Modalit\`a \code{DEPLOYMENT} & Non disponibile: nessun provider. \\
Limite di concorrenza & Nessuna applicazione: senza code non c'\`e chi conti gli
slot. \\
Metriche & Complete per le grandezze del nucleo. \\
\bottomrule
\end{tabular}
\end{adjustbox}
\end{center}

\`E ancora una piattaforma FaaS secondo la definizione operativa del
\cref{cap:02-stato-dell-arte}? Soddisfa l'obbligo di registrazione, quello di
instradamento, e --- in modalit\`a \code{EXTERNAL} --- opera su unit\`a di calcolo
che qualcun altro ha attivato. Non soddisfa attivazione e ritiro. \`E quindi una
piattaforma FaaS \emph{ridotta}, come il \cref{cap:03-requisiti-principi}
prometteva, non una piattaforma rotta: ogni operazione che accetta la porta a
termine correttamente, e le operazioni che non pu\`o svolgere le rifiuta
esplicitamente.

\subsection{\code{501}, non \code{503}}

La scelta del codice di stato per \code{:enqueue} \`e deliberata. \code{503
Service Unavailable} significa ``riprova pi\`u tardi''; \code{501 Not
Implemented} significa ``questa capacit\`a non esiste in questo binario''. Un
client che riceve \code{503} riprova e continuer\`a a fallire; uno che riceve
\code{501} sa che deve cambiare qualcosa.

\subsection{La copertura di test}

La configurazione senza moduli ha una classe di test dedicata,
\code{CoreOnlyApiTest}, che verifica ci\`o che vale \emph{solo} in sua assenza:
che l'accodatore sia il no-op e che l'API asincrona risponda \code{501}. La
classe si auto-esclude quando i moduli sono presenti.

Il costo \`e reale. Poich\'e la selezione modulare \`e risolta a tempo di
configurazione di Gradle e non pu\`o variare all'interno di una singola build,
quella classe richiede una propria invocazione:

\begin{lstlisting}[language=bash]
./gradlew test                                   # tutti i moduli (default: all)
./gradlew test -PcontrolPlaneModules=none        # la configurazione minima
\end{lstlisting}

Due invocazioni della suite per ogni verifica completa. \`E il prezzo diretto
della scelta di selezionare a build time anzich\'e a runtime
(\cref{cap:08-sistema-moduli}), ed \`e utile averlo quantificato: la
modularit\`a a costo zero sul percorso critico non \`e a costo zero sulla catena
di integrazione continua.

\subsection{Perch\'e il default \`e \code{all}}

Ogni invocazione di Gradle priva di selezione esplicita usa \code{all}, test
compresi. La ragione \`e che la suite deve esercitare la configurazione che viene
distribuita. Una suite che girasse per default su una configurazione parziale
verificherebbe un artefatto che nessuno esegue, e le interazioni fra moduli ---
che sono la parte pi\`u fragile di un sistema componibile --- resterebbero
scoperte.

La combinazione delle due scelte, default \code{all} e test dedicato per
\code{none}, copre i due estremi dello spazio delle configurazioni. Le
configurazioni intermedie --- per esempio \code{async-queue} senza
\code{sync-queue} --- non hanno copertura dedicata.

\begin{damisurare}
Lo spazio delle configurazioni modulari ha $2^9 = 512$ punti, di cui alcuni non
validi per la dipendenza fra \code{concurrency-control} e \code{async-queue}. La
suite ne verifica due. Una verifica di costruibilit\`a su tutte le combinazioni
valide \`e realizzabile e non \`e stata eseguita.
\todomis{costruibilit\`a e avviabilit\`a di tutte le combinazioni modulari valide}
\end{damisurare}
